\section{Conclusion}

We introduced a data-free lineage detection score derived from checkpoint-specific structure in residual branch products. Across residual-MLP and GPT-2 benchmarks, the score separated fine-tuned, LoRA-merged, pruned, and quantized descendants from independently trained and distilled models, while distinguishing weight ancestry from behavioral similarity. It remained stable under function-preserving checkpoint laundering and matched the nearest robust baseline at lower latency. Projection-pairing results across six language-model families and a public LLaMA-2 case study support the approach. The method is limited to compatible white-box residual checkpoints and does not establish ownership or direction of descent. It enables provenance auditing for open-weight language-model supply chains.

\section*{Ethical Considerations}

We present a method for detecting shared weight ancestry between neural network checkpoints. Potential applications include model provenance verification and supply-chain auditing. We acknowledge several ethical considerations:

\textbf{Limitations of evidence.} The method produces a similarity score, not a legal determination. Scores consistent with documented relationships do not independently prove provenance, ownership, or wrongdoing. Results should be combined with metadata, release records, and human review before drawing conclusions about model origins.

\textbf{False confidence risk.} Limited calibration data (few independent roots) may produce misleadingly precise-seeming thresholds. Users should understand that formal hypothesis testing requires sufficient exchangeable samples, which may not be available for rare architectures.

\textbf{Privacy considerations.} The method could potentially be used to identify undocumented model relationships, which may conflict with developers' reasonable expectations of anonymity. The white-box access requirement limits covert use, but authorized access does not automatically imply consent to provenance analysis.

\textbf{Dual use.} The same method that enables legitimate supply-chain auditing could be misused for unfounded infringement claims or competitive intelligence. We encourage responsible deployment with appropriate evidentiary standards.

\section*{Limitations}\label{sec:limitations}

\textbf{Scope and architecture constraints.} Our method needs white-box access to both checkpoints, so API-only models cannot be verified. It also only works for residual architectures: plain feedforward networks, RNNs, and state-space models are out of scope. Comparisons are further limited to models with matching depth and hidden dimension, so cross-architecture verification (e.g., LLaMA vs.\ GPT-2) is not supported.

\textbf{Lineage determination.} The lineage score is symmetric, so without external metadata we cannot tell which checkpoint is the ancestor. The score decays monotonically with genealogical distance, but the method compares single pairs and does not reconstruct multi-hop ancestry or family trees. Under partial inheritance (layer grafts, linear merges), the score reflects the fraction of shared blocks but not which blocks were inherited.

\textbf{Signal degradation under transformation.} The fingerprint survives common post-training modifications, but weakens as transformations grow more aggressive. Heavy pruning drops the signal to $\mathcal{L}=0.58$ at 85\% sparsity, and extensive continued pretraining (CodeLlama) lowers it to $\mathcal{L}=0.336$, though both remain detectable above the null. In the evaluated MLP suppression attack, reaching the null threshold ($\mathcal{L} \approx 0.084$) costs +1.5\% utility loss; driving the score reliably below null costs +12\%. This result does not establish robustness to other attacks, architectures, utility measures, or optimization budgets.

\textbf{Residual-stream rotation.} The signature is invariant to hidden-unit permutation and reciprocal rescaling within MLP blocks, but not to orthogonal rotation $Q$ of the residual stream: $M \mapsto QMQ^\top$ preserves trace but transforms the centered remainder, destroying cross-checkpoint similarity. For LayerNorm architectures (GPT-2, BERT), such rotation is not function-preserving because learned element-wise $\gamma$ and $\beta$ parameters do not commute with arbitrary $Q$. For RMSNorm architectures with learned per-dimension scale (LLaMA, Mistral), the same obstruction applies unless $\gamma$ is constant across dimensions, which is not standard practice. An architecture using only unparameterized normalization (e.g., pure RMS without learned scale) would be vulnerable to $Q$-rotation laundering; we are not aware of widely deployed models in this category.

\textbf{Calibration and validation.} Verification requires calibrating a null distribution from independently trained models of the same architecture, and the threshold may need adjustment across architecture families. Our controlled benchmarks demonstrate perfect AUROC discrimination, but the effective sample sizes for formal hypothesis testing are limited: the GPT-2-style benchmark has 8 roots total. The LLaMA-2 case study now includes 7 independent models (OpenLLaMA v1/v2, Amber, Baichuan v1/v2, InternLM, Yi), providing stronger calibration for that architecture family, but other families remain undertested. Real-world validation covers a limited set of public checkpoints, so broader ecosystem coverage is untested. SwiGLU architectures (Qwen2.5, DeepSeek-R1) reach only 80\% block-pairing accuracy and need joint factorization to recover sub-paths reliably.

\section*{Use of AI Assistants}

The authors bear full responsibility for all scientific claims, experimental results, and analysis in this work. During manuscript preparation, we used Claude to support literature searches, experimental design, code development, and the drafting and editing of text. All AI-assisted material was reviewed, verified, and revised by the authors.
